\section{Small connection eigenvalues and mean curvature}
\label{sec:weak-spectrum}

We first identify all small eigenvalues of a real connection operator for
the compact metrics \(H_a\). The argument applies to an arbitrary finite
defect. We then use the normalized reference metric of
Lemma~\ref{lem:normalized-reference} to estimate the mean curvature when
the defect is reduced. All constants in this section concern the fixed
extension, base metric and reference metric.

\subsection{The real connection operator}

Let \(D_a\) be the Chern connection of \(H_a\). On the real bundle of
\(H_a\)-Hermitian trace-free endomorphisms, set
\begin{equation}\label{eq:connection-form}
 \mathfrak q_a(s,t)
 =\operatorname{Re}\int_X\langle D_a''s,D_a''t\rangle\,\dd V
 =\frac12\operatorname{Re}\int_X\langle D_as,D_at\rangle\,\dd V,
 \qquad \|s\|_a^2=\int_X\tr(s^2)\,\dd V.
\end{equation}
The second equality follows from Hermitian adjunction and real
polarization. The associated self-adjoint operator is
\(\mathcal L_a=\frac12D_a^*D_a\). We write its eigenvalues, counted with
multiplicity, as \(\lambda_1(a)\le\lambda_2(a)\le\cdots\).
At a metric with noncentral mean curvature the complex operator
\((D_a'')^*D_a''\) on all endomorphisms need not preserve Hermitian
sections. The real operator in \eqref{eq:connection-form} is the one
used throughout the proof.

Let \(\mathcal P\) be the space of parallel Hermitian trace-free
endomorphisms of the smooth graded connection \(D_*\) in
\eqref{eq:graded-connection}. Then
\begin{equation}\label{eq:weak-dimension}
 k=\dim_{\RR}\mathcal P=
 \begin{cases}
  1,&L\not\simeq M,\\
  3,&L\simeq M.
 \end{cases}
\end{equation}
Indeed a nonzero parallel off-diagonal map is a line-bundle
isomorphism. If the two line bundles are isomorphic, their HE metrics
differ, after identification, by a positive constant; the logarithmic
ratio is harmonic. A normalized isomorphism is therefore parallel.
It supplies two real off-diagonal directions in addition to the
diagonal trace-free direction.

\begin{proposition}\label{prop:weak-spectrum}
For the family \eqref{eq:compact-metric} with any smooth reference
metric of the fixed determinant, all \(\lambda_j(a)\) are positive.
There are \(a_0,c>0\) such that
\begin{equation}\label{eq:complete-cluster}
 \lambda_k(a)\longrightarrow0,\qquad
 \lambda_{k+1}(a)\ge c\quad(0<a\le a_0).
\end{equation}
If \(W\) is reduced, then \(\lambda_k(a)\le Ca\).
The small spectral space converges, on every compact subset of
\(X\setminus\Supp W\), to \(\mathcal P\) under the unitary
identifications specified below. In particular, it contains every
normalized mode whose connection energy tends to zero.
\end{proposition}

\begin{proof}
Put \(U=X\setminus\Sigma\), where \(\Sigma=\Supp W\), and let \(\tau\)
be the \(H_0\)-orthogonal right inverse of the quotient map.
The splitting \((\iota,\tau)\) has metric
\(\diag(u_ah_L,u_a^{-1}h_M)\). Consequently
\begin{equation}\label{eq:graded-unitary-map}
 F_a(l,m)=u_a^{-1/2}\iota(l)+u_a^{1/2}\tau(m)
\end{equation}
is a unitary isomorphism from \((L\oplus M,h_L\oplus h_M)\) to
\((V,H_a)\) over \(U\). Write \(\db\tau=\iota\beta\), with the
fixed line metric on \(\operatorname{Hom}(M,L)\) understood.
The diagonal connection differences in this unitary splitting are
opposite copies of
\(\frac12(\partial-\db)\log u_a\); the off-diagonal blocks have
norm \(u_a|\beta|_0\). Thus
\begin{equation}\label{eq:graded-connection-error}
 |F_a^*D_a-D_*|
 \le C\bigl(|\dd\log u_a|+u_a|\beta|_0\bigr).
\end{equation}
Since \(\log u_a=\log a-\log(1+a/u_0)\), this difference and all
its fixed derivatives are \(O_K(a)\) on each \(K\Subset U\).

The scalar Sobolev inequality on the fixed four-dimensional base,
together with Kato's inequality, gives
\begin{align}
 \|s\|_{L^4(a)}^2
 &\le C\bigl(2\mathfrak q_a(s,s)+\|s\|_a^2\bigr),
                                                   \label{eq:uniform-sobolev}\\
 \int_{B_\rho(\Sigma)}|s|_a^2\,\dd V
 &\le C\rho^2\bigl(2\mathfrak q_a(s,s)+\|s\|_a^2\bigr).
                                                   \label{eq:no-point-mass}
\end{align}
The second estimate uses the \(O(\rho^4)\) volume of the finitely
many balls. Both are intrinsic estimates on the original smooth
bundle and require no curvature bound.

Suppose \(a_n\downarrow0\), \(\|s_n\|_{a_n}=1\), and
\(\mathfrak q_{a_n}(s_n,s_n)\to0\).
Transport by \(F_{a_n}\) on fixed compact subsets of \(U\).
The convergence of the connection coefficients and the covariant
derivative bound give local \(H^1\) bounds. Rellich compactness,
followed by a diagonal subsequence, gives strong local \(L^2\)
convergence to a section \(s_*\) satisfying \(D_*s_*=0\).
It is smooth by elliptic regularity. Its holomorphic coefficients
extend over the isolated points by Hartogs' theorem, and the
extension is parallel. Estimate~\eqref{eq:no-point-mass}, first
with fixed \(\rho\) and then with \(\rho\downarrow0\), shows that
the limit has norm one. Applying the same argument to sums of
finitely many sequences preserves their real inner products.

To obtain test sections, choose a cutoff \(\chi_\rho\) equal to
zero on \(B_\rho(\Sigma)\), equal to one outside \(B_{2\rho}(\Sigma)\),
and satisfying \(|\dd\chi_\rho|\le C/\rho\). For \(b\in\mathcal P\),
extend
\[
 T_{a,\rho}b=\chi_\rho F_abF_a^{-1}
\]
by zero across the points. Its mass is exactly
\(\int\chi_\rho^2|b|^2\,\dd V\). For fixed \(\rho\), its energy
tends as \(a\downarrow0\) to
\(\frac12\int|\dd\chi_\rho|^2|b|^2\,\dd V\le C\rho^2|b|^2\).
These statements are uniform on the unit sphere of \(\mathcal P\).
Min--max, first with \(\rho\) fixed and then with
\(\rho\downarrow0\), proves \(\lambda_k(a)\to0\).
If the complementary bound in \eqref{eq:complete-cluster} failed,
there would be \(k+1\) orthonormal eigensections with vanishing
energy. The preceding compactness argument would give \(k+1\)
orthonormal elements of \(\mathcal P\), a contradiction.

For positivity at a fixed \(a\), a zero mode would be a nonzero
parallel Hermitian trace-free endomorphism. Its two eigenlines
would split \(V=A\oplus B\) holomorphically. Projecting \(\iota\)
to the factors gives sections of \(A\otimes L^{-1}\) and
\(B\otimes L^{-1}\). If both are nonzero, their zero divisors
have nonnegative Kähler degrees whose sum is
\[
 \deg A+\deg B-2\deg L=\deg V-2\deg L=0.
\]
Both sections are then nowhere zero, contrary to the nonempty
zero scheme of \(\iota\). If only one projection is nonzero,
its zero set would be a nonempty finite set; a line-bundle
section on a smooth surface has either a divisor of zeros or
no zeros. This is again impossible. Thus \(\ker\mathcal L_a=0\).
This argument does not assert simplicity of \(V\); a holomorphic
nilpotent endomorphism in the repeated-factor case is not
Hermitian and parallel.

If \(W\) is reduced, in each fixed coordinate ball
\[
 u_0\asymp r^2,\qquad |\dd u_0|\le Cr,\qquad
 |\beta|_0\le Cr^{-3}.
\]
For the last bound, the numerator in
\(\tau=H_0^{-1}q^*/(qH_0^{-1}q^*)\) is \(O(r)\) and the
denominator is comparable to \(r^2\). Hence
\(\db\tau=O(r^{-2})\), and division by the inclusion,
whose norm is comparable to \(r\), gives the assertion.
For \(r\ge\sqrt a\), \eqref{eq:graded-connection-error} is
therefore bounded by \(Ca/r^3\). Choose the cutoff at radius
\(\sqrt a\). Its derivative contributes \(O(a)\) to the energy,
and the connection term contributes at most
\[
 Ca^2\int_{\sqrt a}^{R}r^{-6}r^3\,\dd r\le Ca.
\]
The fixed exterior contributes \(O(a^2)\), whereas the omitted
mass is \(O(a^2)\). Min--max proves \(\lambda_k(a)\le Ca\).

For precision about the subspaces, let \(\Pi_a\) be the orthogonal
projection onto the first \(k\) eigensections and let \(T_a\)
be the just constructed cutoff map. With the full graded
\(L^2\) inner product on \(\mathcal P\), the complementary gap gives
\begin{equation}\label{eq:spectral-identification}
 \|(I-\Pi_a)T_ab\|_a^2\le Ca\|b\|_*^2,\qquad
 R_a^*R_a=I+O(a),\qquad R_a=\Pi_aT_a.
\end{equation}
Thus \(R_a\) is an isomorphism onto the complete small spectral
space. On fixed compact subsets the transported \(T_ab\) equals
\(b\), so \(R_ab\to b\) in local \(L^2\). Expand \(R_ab\) in a
normalized low eigenbasis. Its coefficients are bounded, and the
transported elliptic operators have uniformly smooth coefficients
there. Interior estimates give convergence in every fixed
\(C^j\) norm. For a general finite \(W\), choose a sequence
\(\rho_j\downarrow0\) and then thresholds \(a_j\downarrow0\)
so that the fixed-\(\rho_j\) connection error is at most
\(\rho_j^2\) whenever \(a\le a_j\). Using those cutoffs successively
gives mass \(I+o(1)\) and energy \(o(1)\), and the same proof
applies. No parameter regularity of that last cutoff choice is needed.
\end{proof}

\subsection{The normalized mean-curvature residual}

For the remainder of this section assume \(W\) is reduced and use
Lemma~\ref{lem:normalized-reference}. Choose a smooth positive weight
\(\rho_a\) comparable to \(\sqrt{a+|z|^2}\) on the fixed balls and
comparable to one on their complement. For example one can take the
square root of \(a\) plus a fixed smooth function equal to \(|z|^2\)
near each point and positive elsewhere. Define
\[
 \Phi_a=i\Lambda_\omega F_{H_a}-\kappa I,\qquad
 \ell(a)=1+|\log a|.
\]

\begin{proposition}\label{prop:residual}
For each fixed nonnegative integer \(j\),
\begin{equation}\label{eq:weighted-residual}
 |D_a^j\Phi_a|_{H_a}\le C_j a\rho_a^{-2-j}.
\end{equation}
Consequently
\begin{equation}\label{eq:residual-lp}
 \|\Phi_a\|_p\le
 \begin{cases}
  C_pa,&1\le p<2,\\
  Ca\sqrt{\ell(a)},&p=2,\\
  C_pa^{2/p},&2<p<\infty,
 \end{cases}
 \qquad \|\Phi_a\|_\infty\le C.
\end{equation}
The smooth parameter direction
\begin{equation}\label{eq:parameter-direction}
 k_a=H_a^{-1}a\partial_aH_a
     =\frac{2S-u_0I}{u_0+a}
\end{equation}
satisfies
\begin{equation}\label{eq:parameter-pairing}
 \|k_a\|_a^2\longrightarrow2v_0,\qquad
 \langle\Phi_a,k_a\rangle_a
 =\frac{2\mathcal E_a}{a}=4\pi^2\sigma a+o(a).
\end{equation}
These assertions allow \(L\simeq M\).
\end{proposition}

\begin{proof}
In the normalized frames near \(p\), put \(s=|z|^2\) and
\(c=\lambda_p\sqrt{h_M/h_L}\). Then
\[
 u_0=cs,\quad u_a=\frac{acs}{a+cs},\quad
 c=\lambda_p+O(r^2),\quad \dd c=O(r),
 \quad g=g_{\mathrm E}+O(r^2).
\]
The orthogonal quotient lift and its \(\db\) derivative are
\begin{equation}\label{eq:normalized-splitting}
 \tau=\lambda_p^{-1/2}\frac{(\bar z_2,\bar z_1)^{\mathsf t}}s,
 \qquad
 \beta=\frac{\bar z_1\dd\bar z_2-\bar z_2\dd\bar z_1}
                  {\lambda_p s^2}.
\end{equation}
Since the line weight on \(\operatorname{Hom}(M,L)\) is
\((\lambda_p/c)^2\) and \(|\dd z_j|^2=2\),
\begin{equation}\label{eq:beta-leading}
 |\beta|_0^2=\frac{2}{c^2s^3}+O(s^{-2}).
\end{equation}
Write \(\Delta=-\Delta_0\) for the nonpositive real scalar
Laplacian. In the unitary splitting the diagonal residual is
\begin{equation}\label{eq:diagonal-residual}
 (\Phi_a)_{11}=u_a^2|\beta|_0^2-\frac12\Delta\log u_a,
 \qquad(\Phi_a)_{22}=-(\Phi_a)_{11}.
\end{equation}
For example, the first identity follows directly by substituting
\eqref{eq:full-inclusion-energy} into
\(\Delta u_a=2e_a-2u_a(\Phi_a)_{11}\). The trace identity follows
from the fixed HE determinant.

For \(v=cs\) one has
\(\Delta v=8c+O(s)\) and
\(|\dd v|^2=4c^2s+O(s^2)\). The derivatives of
\(f(v)=\log a+\log v-\log(v+a)\) are
\[
 f'(v)=\frac{a}{v(v+a)},\qquad
 f''(v)=-\frac{a(2v+a)}{v^2(v+a)^2}.
\]
The chain rule therefore gives
\begin{align}
 \Delta\log u_a
 &=\frac{4a^2}{s(a+cs)^2}
       +O\left(\frac{a}{a+s}\right),\\
 u_a^2|\beta|_0^2
 &=\frac{2a^2}{s(a+cs)^2}
       +O\left(\frac{a^2}{(a+s)^2}\right).
\end{align}
For the first error, multiply \(O(s)\) by
\(|f'|\le Ca/[s(s+a)]\), and \(O(s^2)\) by
\(|f''|\le Ca/[s^2(s+a)]\). The leading terms cancel in
\eqref{eq:diagonal-residual}, leaving \(Ca/(a+s)\).

The upper off-diagonal curvature block in the unnormalized
splitting is
\[
 D'_0\beta+2\partial\log u_a\wedge\beta.
\]
After unitary normalization its contraction has norm at most
\(Cu_a|\Lambda(D'_0\beta+2\partial\log u_a\wedge\beta)|_0\).
Direct differentiation in the flat metric gives
\[
 \Lambda_{\mathrm E}\partial\beta=0,\qquad
 \Lambda_{\mathrm E}(\partial s\wedge\beta)=0.
\]
The contraction error is \(O(s)\), the fixed line connection is
\(O(r)\), and \(\partial\beta=O(r^{-4})\). It follows that
\(|\Lambda D'_0\beta|_0\le Cs^{-1}\); multiplication by
\(u_a\le Cas/(a+s)\) gives the required bound. Finally
\[
 \partial\log u_a
 =\frac{a}{s(a+cs)}\partial s+
   \frac{a}{c(a+cs)}\partial c.
\]
The flat contraction of the first term with \(\beta\) vanishes.
The base-metric error and the second term each have norm at most
\(Ca/[s(a+s)]\) before multiplication by \(u_a\). Their
normalized contributions are at most \(Ca^2/(a+s)^2\).
Hermitian symmetry treats the other off-diagonal block. This
proves \eqref{eq:weighted-residual} for \(j=0\).

On \(r\ge\sqrt a\), differentiating these formulas and their
cancellations costs at most \(C_jr^{-j}\). The connection in
the normalized splitting is \(O(r^{-1})\), with derivatives
of order \(j\) bounded by \(C_jr^{-1-j}\). Conversion to
covariant derivatives preserves the estimate
\(C_jar^{-2-j}\). At the center use the smooth holomorphic
frame instead. Under \(z=\sqrt a\,y\), on a fixed \(y\)-ball,
\[
 H_a(\sqrt a\,y)=H^{\mathrm{flat}}_{\lambda_p}(y)+O_{C^j}(a),
 \qquad g(\sqrt a\,y)=g_{\mathrm E}+O_{C^j}(a).
\]
The leading metric is the local HYM metric of
Lemma~\ref{lem:local-metric}, with a fixed dilation and sign
change. Its metric and inverse are uniformly bounded on
the fixed ball. The physical curvature contraction has a
factor \(a^{-1}\), while its leading term is zero and its
error is \(O_{C^j}(a)\). Thus the physical \(\Phi_a\) and its
\(y\)-covariant derivatives are bounded. Returning to physical
derivatives costs \(a^{-j/2}\), which is precisely
\eqref{eq:weighted-residual} on \(r\le2\sqrt a\).
On the fixed exterior, \eqref{eq:graded-connection-error}
gives \(O(a)\) in every fixed derivative. This proves the
global estimate. Integrating
\[
 Ca^p\int_0^R\frac{r^3}{(a+r^2)^p}\,\dd r
\]
proves \eqref{eq:residual-lp}, including the logarithm at \(p=2\).

Differentiating the two eigenvalues \(w,w^{-1}\) in
\eqref{eq:compact-local-expression}, with \(P\) independent
of \(a\), gives
\[
 k_a=\frac{u_0}{u_0+a}(2P-I)=\frac{2S-u_0I}{u_0+a}.
\]
The second expression is smooth across the points. Its
eigenvalues are opposite, bounded in absolute value by one,
and tend to \(\pm1\) off \(\Sigma\). Bounded convergence gives
the mass limit in \eqref{eq:parameter-pairing}. Since \(\Phi_a\)
is trace-free,
\[
 \tr(\Phi_ak_a)=\frac2a u_a(\Phi_a)_{11}.
\]
The integrated Bochner identity for the smooth global section
\(\iota\) says that the integral of \(u_a(\Phi_a)_{11}\) is
\(\mathcal E_a\). There is no puncture boundary term. The
balanced static energy \eqref{eq:balanced-static} now proves
the exact pairing and its asymptotic.
\end{proof}

The residual is small in each finite \(L^p\) space, but the estimate
does not tend to zero in \(L^\infty\). A correction is therefore
needed before \(H_a\) can provide a sufficiently accurate comparison
for the flow. The complete dimension in \eqref{eq:weak-dimension}
will also be retained in that correction.
